% !TeX program = xelatex
\documentclass[11pt]{article}

\usepackage[a4paper, margin=1in]{geometry}
\usepackage{amsmath, amssymb, amsthm}
\usepackage{stmaryrd}
\usepackage{kotex}
\usepackage{listings}
\usepackage{listingsutf8}
\usepackage{xcolor}
\usepackage{hyperref}

% --- Line-breaking inside \texttt{...} ---------------------------------------
% \texttt{} 안의 긴 식별자가 줄 끝에서 끊기지 않아 오른쪽 여백을
% 벗어나는 문제를 피하기 위해, sloppy + emergencystretch로 단락 단위의
% 신축 여유를 크게 주고, \btt{...}는 underscore마다 line-break을 허용하는
% texttt 변종으로 정의한다.  본문의 길게 이어지는 식별자에는 \btt를
% 사용한다.
\sloppy
\setlength{\emergencystretch}{3em}

% \btt{a_b_c} 형태에서 underscore마다 break point를 허용
\newcommand{\btt}[1]{\texttt{\detokenize{#1}}}
% -----------------------------------------------------------------------------

\lstdefinelanguage{lean}{
  morekeywords={theorem, lemma, def, by, fun, let, in, do, match, with,
    if, then, else, return, import, open, namespace, end, structure,
    inductive, partial, private, syntax, elab, elab_rules, macro,
    macro_rules, tactic, conv, command, example, have, show, intro,
    apply, simp, rw, ring, all_goals, repeat},
  sensitive=true,
  morecomment=[l]{--},
  morecomment=[s]{/-}{-/},
  morestring=[b]"
}

\lstset{
  inputencoding=utf8,
  language=lean,
  basicstyle=\ttfamily\small,
  keywordstyle=\color{blue},
  commentstyle=\color{gray}\itshape,
  stringstyle=\color{red!60!black},
  showstringspaces=false,
  breaklines=true,
  columns=fullflexible,
  frame=single,
  framesep=4pt,
  xleftmargin=4pt,
  xrightmargin=4pt
}

\title{Domain-Specific Tactics in LeanFlagAlgebras}
\author{}
\date{}

\begin{document}
\maketitle

\section{Introduction}

The aim of this project is to mechanically verify extremal graph density
inequalities. In doing so we run into many proof patterns that recur
across the development; writing them out by hand would mean repeating
the same \texttt{rw}/\texttt{simp}/\texttt{ring} sequences for tens to
hundreds of lines at a time. To avoid that, we implemented a number of
domain-specific tactics using Lean 4's metaprogramming facilities.

\bigskip\noindent
The tactics fall into three groups, organized below so that each later
group builds on the names and lemmas introduced by the earlier ones:

\begin{itemize}
  \item \textbf{Meta-commands.} These read flag definitions, density
    values, and product expansions from external JSON files and
    register the corresponding definitions and lemmas in bulk. Their
    job is to set up the named definitions/theorems that later tactics
    will look up.
    \texttt{load\_empty\_typed\_flags},
    \texttt{load\_flags},
    \texttt{load\_flag\_pair\_density\_theorems},
    \texttt{load\_forbid\_density\_theorems},
    \texttt{load\_forbid\_mul\_theorems}.
  \item \textbf{Flag-algebra-specific tactics.} These use the naming
    convention set up by the meta-commands to unfold and rearrange
    flag algebra expressions in actual proofs.
    \texttt{expand\_one\_at},
    \texttt{reduce\_flagmul},
    \texttt{reduce\_downward\_flagmul},
    \texttt{prove\_flag\_expand\_with\_forbidden\_flag},
    \texttt{prove\_flag\_mul\_with\_forbidden\_flag},
    \texttt{flag\_nonneg}.
  \item \textbf{Algebraic-normalization tactics.} Used at the closing
    step to sort summands and discharge equalities up to add-AC.
    The \texttt{sort}/\texttt{ac\_sort} family.
\end{itemize}

\section{Meta-commands}

\subsection{\texttt{load\_empty\_typed\_flags "<file>.json"}}

Defined in \texttt{LeanFlagAlgebras/Flags/FlagLoader.lean} (with a
universe-tuned variant of the same name in
\texttt{Flags/UnivEffi.lean}). The command reads a JSON file
\texttt{graphs\_n.json} that enumerates the empty-typed flags of size
$n$ (i.e.\ unlabeled flags whose type is empty) and registers, in one
go, the following:

\begin{itemize}
  \item For each flag $i$, the chain
    \texttt{Sym2Graph\_n\_0\_0\_i} (the raw \texttt{Sym2Graph n} built
    from the JSON edge list; \texttt{edges\_valid} closed by
    \texttt{decide})
    \texttt{ → Sym2Flag\_n\_0\_0\_i} (its isomorphism class, via
    \texttt{Quotient.mk})
    \texttt{ → Flag\_n\_0\_0\_i} (the bridge into the general flag
    category)
    \texttt{ → noncomputable def FlagAlgebra\_n\_0\_0\_i}
    $= \llbracket \mathrm{unitVector}\;\langle n,\mathrm{Flag}\_n\_0\_0\_i\rangle\rrbracket$.
  \item The finsets \texttt{Sym2FlagSet\_n\_0\_0} and
    \texttt{flagSet\_n\_0\_0} of all such flags, together with the
    three lemmas
    \texttt{Sym2FlagSet\_n\_0\_0\_eq\_univ},
    \texttt{flagSet\_n\_0\_0\_val\_eq} (whose \texttt{.val} matches
    the list we wrote down), and
    \texttt{flagSet\_n\_0\_0\_eq\_univ}, all discharged by
    \texttt{native\_decide}.
\end{itemize}

The \texttt{graphs\_n.json} files are not written by hand. A short
Python script, \texttt{LeanFlagAlgebras/Flags/generate\_graphs.py},
takes the target size $n$ and emits the canonical edge-list
enumeration of all isomorphism classes of graphs on $n$ vertices,
using NetworkX's graph atlas together with a lexicographic-minimal
relabeling pass to pick a canonical representative for each class.
Supporting a new size is therefore one script run plus one
\texttt{load\_empty\_typed\_flags} call. The project loads sizes
$0,\ldots,5$ together in \texttt{LeanFlagAlgebras/Flags/FlagDef.lean}:

\begin{lstlisting}
import LeanFlagAlgebras.Flags.FlagLoader

load_empty_typed_flags "LeanFlagAlgebras/Flags/Graphs/graphs_0.json"
...
load_empty_typed_flags "LeanFlagAlgebras/Flags/Graphs/graphs_5.json"
\end{lstlisting}

After this, names like \texttt{Flag\_5\_0\_0\_17} are available
project-wide, and tactics such as \texttt{expand\_one\_at 5} pick up
the matching \texttt{flagSet\_5\_0\_0\_eq\_univ} and
\texttt{flagSet\_5\_0\_0\_val\_eq} produced by this call.

Internally the command reads the JSON via
\texttt{IO.FS.readFile + Json.parse} and extracts $n$ from the
filename via \texttt{parseNFromGraphsPath}, forcing
\texttt{graphs\_5.json} to correspond exactly to size~$5$. Edge
validity is checked by \texttt{decide} and the completeness statement
by \texttt{native\_decide}, so both the soundness and completeness of
the externally generated enumeration are re-verified inside Lean. The
guard \texttt{env.contains} skips names that are already defined,
allowing repeated or nested imports without conflict.

\subsection{\texttt{load\_flags "<file>.json"}}

The labeled-flag counterpart of \texttt{load\_empty\_typed\_flags},
defined in the same files. \texttt{load\_flags} reads a JSON file
listing all flags of size $n$ on a fixed type of size $k$ with type
index $m$, and registers:

\begin{itemize}
  \item The type graph and its bridge:
    \texttt{Sym2FlagType\_k\_m}, \texttt{FlagType\_k\_m}.
  \item For each flag $i$, the chain
    \texttt{Sym2LabeledGraph\_n\_k\_m\_i} (carries both the edge set
    and the type embedding $\mathrm{Fin}\;k \hookrightarrow
    \mathrm{Fin}\;n$, synthesized automatically)
    \texttt{ → Sym2Flag\_n\_k\_m\_i}
    \texttt{ → Flag\_n\_k\_m\_i}
    \texttt{ → noncomputable def FlagAlgebra\_n\_k\_m\_i}
    (analogous to the empty-typed case).
  \item Two \texttt{@[simp]} lemmas connecting back to the
    empty-typed flags registered by \texttt{load\_empty\_typed\_flags}:
    \begin{itemize}
      \item \texttt{unlabel\_n\_k\_m\_i} ― forgetting the labeling
        gives the underlying empty-typed flag,
        \texttt{FlagAlgebras.unlabel Flag\_n\_k\_m\_i =
        Flag\_n\_0\_0\_{\underline i}}.
      \item \texttt{downward\_n\_k\_m\_i} ―
        $\llbracket \mathrm{FlagAlgebra}\_n\_k\_m\_i \rrbracket_0
        = c_{n,k,m,i}\cdot
        \mathrm{FlagAlgebra}\_n\_0\_0\_{\underline i}$, where the
        coefficient $c$ (the ``downward normalizing factor'') is
        stored as the \texttt{"num/den"} string in the JSON.
    \end{itemize}
  \item The finsets \texttt{sym2FlagSet\_n\_k\_m} and
    \texttt{flagSet\_n\_k\_m} together with the three completeness
    lemmas
    \texttt{sym2FlagSet\_n\_k\_m\_eq\_univ},
    \texttt{flagSet\_n\_k\_m\_val\_eq},
    \texttt{flagSet\_n\_k\_m\_eq\_univ}.
\end{itemize}

As with empty-typed flags, the JSON files come from a Python script:
\texttt{LeanFlagAlgebras/Flags/generate\_flags.py} takes the three
integers $n$, $k$, $m$ and emits \texttt{flags\_n\_k\_m.json},
including the type embedding and the downward normalizing factor for
each flag. The project loads all required $(n,k,m)$ combinations
together in \texttt{LeanFlagAlgebras/Flags/FlagDef.lean}:

\begin{lstlisting}
load_flags "LeanFlagAlgebras/Flags/Flags/flags_1_1_0.json"
load_flags "LeanFlagAlgebras/Flags/Flags/flags_2_1_0.json"
...
load_flags "LeanFlagAlgebras/Flags/Flags/flags_5_3_2.json"
\end{lstlisting}

The internals follow the same pattern as
\texttt{load\_empty\_typed\_flags}, with two extra ingredients. The
type embedding $\mathrm{Fin}\;k \hookrightarrow \mathrm{Fin}\;n$ is
synthesized from the \texttt{type\_indices} array by
\texttt{mkTypeIndexNatExpr} as an if-else ladder, and both injectivity
and adjacency-preservation are discharged by
\texttt{fin\_cases <;> decide}. The downward coefficient $c$ is parsed
from its \texttt{"num/den"} form into a $\mathbb{Q}$ term and the
theorem body is closed by
\texttt{downwardNormalizingFactor\_eq} together with
\texttt{native\_decide}, so the numerical value is re-derived from the
definition rather than trusted; a miscomputation by the script would
surface as a compile error. The \texttt{unlabel\_*} and
\texttt{downward\_*} lemmas are registered with \texttt{@[simp]},
which lets downstream \texttt{simp} calls handle unlabeling and
downward conversion automatically.

\subsection{\texttt{load\_flag\_pair\_density\_theorems "<file>.json"}}

This command, defined in
\texttt{LeanFlagAlgebras/Flags/Densities/DensityLoader.lean}, reads a
JSON file produced by \texttt{calculate\_densities.py} containing
pre-computed pair densities $d_2(F_{p_1}, F_{p_2}; F_h)$ for two
pattern flags and a host flag, and registers them as \texttt{@[simp]}
lemmas. For each row \texttt{[p1, p2, h, "num/den"]} in the JSON it
produces:

\begin{itemize}
  \item \texttt{@[simp] theorem flagDensity₂\_\allowbreak
    Flag\_<patTag>\_p1\_\allowbreak
    Flag\_<patTag>\_p2\_\allowbreak
    Flag\_<hostTag>\_h} ― states
    $\mathrm{flagDensity}_2\,F_{p_1}\,F_{p_2}\,F_h
       = \tfrac{\mathrm{num}}{\mathrm{den}}$
    and is closed by
    \texttt{flagDensity₂\_eq\_sym2FlagDensity₂} followed by
    \texttt{native\_decide}.
\end{itemize}

\texttt{API/C4TuranAPI.lean} and \texttt{ErdosPentagon/FlagMul.lean}
call this for each relevant $(\text{host},\text{pattern})$ pair, e.g.\
\texttt{density\_5\_3\_1\_from\_4\_3\_1.json}. One such JSON file
typically contains hundreds to thousands of density values; once
registered as \texttt{@[simp]}, downstream \texttt{simp} calls reduce
any of these density expressions to a rational constant in one step.

\subsection{\texttt{load\_forbid\_density\_theorems "<file>.json"}}

This command, also defined in
\texttt{LeanFlagAlgebras/Flags/Densities/DensityLoader.lean}, reads a
JSON file produced by \texttt{gen\_free\_indices.py} which lists the
indices of size-$n$ empty-typed flags that are $H$-free for a fixed
forbidden subgraph $H$ ($K_3$, $K_4$, $\ldots$). For each flag $F_i$ it
registers one of:

\begin{itemize}
  \item if $F_i$ is $H$-free:
    \texttt{@[simp] theorem flagDensity1\_H\_Flag\_n\_0\_0\_i\_eq\_zero}
    ― states \texttt{flagDensity₁ H.toFinFlag.2 Flag\_n\_0\_0\_i = 0};
  \item otherwise:
    \texttt{@[simp] theorem flagDensity1\_H\_Flag\_n\_0\_0\_i\_ne\_zero}
    ― states \texttt{¬ flagDensity₁ H.toFinFlag.2 Flag\_n\_0\_0\_i = 0}.
\end{itemize}

Each theorem unfolds the forbidden-graph definition, rewrites with
\texttt{flagDensity₁\_eq\_sym2EmptyTypeFlagDensity₁}, and closes by
\texttt{native\_decide}.

\texttt{API/C4TuranAPI.lean}, \texttt{ErdosPentagon/FlagMul.lean},
and \texttt{API/K4freeP4.lean} load the appropriate JSON (e.g.\
\texttt{graphs\_5\_K3\_free\_indices.json}). The ``is $F_i$
$H$-free?'' check that recurs throughout the forbid-conditioned
proofs is thereby precomputed and handed to \texttt{simp}, so
downstream tactics never need to redo the classification.

\subsection{\texttt{load\_forbid\_mul\_theorems "<file>.json"}}

This command, defined in
\texttt{LeanFlagAlgebras/Flags/Densities/MulLoader.lean}, consumes the
\emph{same} JSON file as \texttt{load\_flag\_pair\_density\_theorems}
(produced by \texttt{calculate\_densities.py}) together with the
free-index data registered by \texttt{load\_forbid\_density\_theorems},
and produces ``product expansion'' theorems. For each pair of pattern
free indices $(i,j)$ with $i \le j$ it registers:

\begin{itemize}
  \item \texttt{theorem flagMul\_FlagAlgebra\_<patTag>\_i\_FlagAlgebra\_<patTag>\_j}
    ― states
    \[
       F_i \cdot F_j \;=_{H\text{-forbid}}\;
       \sum_{h \in \text{host\_free}} c_{i,j,h}\;
       \mathrm{FlagAlgebra}\_{<\!\!\mathrm{hostTag}\!\!>}\_h,
    \]
    proved by rewriting with
    \texttt{unitVector\_quot\_mul\_forbidEq\_sum} and expanding the sum
    via \texttt{flagSet\_<hostTag>\_eq\_univ} and
    \texttt{flagSet\_<hostTag>\_val\_eq}, so that forbidden host flags
    are eliminated by \texttt{simp}. (If the original call had $i > j$,
    the proof first applies \texttt{rw [mul\_comm]}.)
\end{itemize}

In practice this command is called in tandem with
\texttt{load\_flag\_pair\_density\_theorems} on the same density JSON,
e.g.\ \texttt{density\_5\_3\_2\_from\_4\_3\_2.json}. The generated
names \texttt{flagMul\_<A>\_<B>} match exactly those later synthesized
by \texttt{reduce\_flagmul} and \texttt{reduce\_downward\_flagmul}, so
the producer and the consumer share a single naming convention. What
would otherwise be tens to hundreds of ``$F_i \cdot F_j = \ldots$''
lemmas per host size, each with a hand-computed right-hand side, is
reduced to one JSON file and one command.

\section{Flag-algebra-specific tactics}

These tactics use the names registered by the meta-commands
(\texttt{Flag\_*}, \texttt{FlagAlgebra\_*}, \texttt{flagSet\_*\_eq\_univ},
\texttt{flagMul\_*}) inside actual proofs. They are presented in the
order they typically appear in a proof: expand $\to$ resolve products
$\to$ use the forbidden assumption $\to$ close by non-negativity.

\subsection{\texttt{expand\_one\_at n}}

Defined in \texttt{LeanFlagAlgebras/API/Basic.lean}. Under a forbid
condition \texttt{F\_forbid}, the tactic rewrites the unit
$1 \in \mathrm{FlagAlgebra}$ as the explicit weighted sum of
empty-typed flags of size $n$ that do not contain \texttt{F\_forbid} as
a subgraph. The underlying definition is
\texttt{forbidExpand\_one F\_forbid n}, also in
\texttt{API/Basic.lean}; \texttt{expand\_one\_at n} unfolds it and
spells out the sum using the \texttt{flagSet\_n\_0\_0\_*} lemmas
registered by \texttt{load\_empty\_typed\_flags}.

It is the first step in all four main API proofs
(\texttt{MantelTheoremAPI}, \texttt{C4TuranAPI}, \texttt{K4freeP4},
\texttt{ErdosPentagonAPI}); the size argument changes from $3$ to $5$
across them but the call site looks the same:

\begin{lstlisting}
expand_one_at 3   -- or 4, or 5 in the other proofs
\end{lstlisting}

The tactic synthesizes the required theorem names
(\texttt{flagSet\_\$n\_0\_0\_eq\_univ},
\texttt{flagSet\_\$n\_0\_0\_val\_eq}) from \texttt{n} and rewrites
with them, so the same call works at every size.

\subsection{\texttt{reduce\_flagmul}}

Defined in \texttt{LeanFlagAlgebras/ErdosPentagon/Lemmas.lean}. The
tactic handles equality goals of the form
\[
   c_1 (A_1 \cdot B_1) + c_2 (A_2 \cdot B_2) + \cdots + c_n (A_n \cdot B_n)
   \;=_{F\text{-forbid}}\; \mathrm{RHS},
\]
i.e.\ a linear combination of flag-algebra products on the left.
Starting from the head, it picks the leftmost product $A_i \cdot B_i$,
looks up the matching \texttt{flagMul\_<A>\_<B>} theorem registered by
\texttt{load\_forbid\_mul\_theorems} (which expands $A_i \cdot B_i$ as
a linear combination of base flags
$\sum_j d_{ij}\,F_j$), uses it to replace the head term, and moves
$c_i\,\sum_j d_{ij}\,F_j$ onto the right-hand side. The next product
becomes the new head, and the process repeats until the left-hand
side is exhausted. The net effect is that every product on the left
has been resolved and the difference relative to \texttt{RHS} now
sits entirely on the right.

A representative use is the lemma
\texttt{flagQuadraticForm\_P\_v0\_forbidEq} in
\texttt{ErdosPentagon/Lemmas.lean}: $v_0$ is an $8$-dimensional vector,
so unfolding \texttt{flagQuadraticForm} yields $8^2 = 64$ product
terms, all dispatched by a single \texttt{reduce\_flagmul} call.

The tactic is essentially a fixed-point loop (fuel $256$): it inspects
the head term as an \texttt{Expr}, extracts the
\texttt{FlagAlgebra\_*}/\texttt{Flag\_*} constants from the two
factors, synthesizes the name \texttt{flagMul\_<A>\_<B>} (trying both
factor orderings, and both the local namespace and
\texttt{ErdosPentagon}), and rewrites with the first match.

\subsection{\texttt{reduce\_downward\_flagmul}}

Defined in \texttt{LeanFlagAlgebras/API/ReduceFlagMul.lean}. The
tactic handles \texttt{forbidLE} goals whose left-hand side is a sum
in which the products on the upper type are wrapped in
\texttt{downward}. Concretely, the goal has the form
\[
   T_1 + T_2 + \cdots + T_n \;\le_{F\text{-forbid}}\; \mathrm{RHS},
\]
where each $T_i$ is one of two shapes:
\begin{enumerate}
  \item \texttt{downward(}$c_i \cdot (A_i \cdot B_i)$\texttt{)} ―
    a product $A_i \cdot B_i$ of two unit flags taken on the upper
    type and then pushed down to the base ($\emptyset_t$) type. The
    product itself cannot be expanded into base-type flags directly:
    it must be expanded on the upper type and only then pushed down.
  \item A single flag term (e.g.\ $\mathrm{FlagAlgebra}\_n\_0\_0\_k$)
    that already lives on the base type and just needs to be moved to
    the right.
\end{enumerate}

For each leading term $T$ the tactic acts as follows:

\begin{itemize}
  \item If $T = \texttt{downward}(c \cdot (A \cdot B))$: it looks up the
    \texttt{flagMul\_<A>\_<B>} lemma registered by
    \texttt{load\_forbid\_mul\_theorems}, which states
    $A \cdot B =_{F\text{-forbid}} \sum_j d_j\,F_j$. Combining this
    with the coefficient $c$ via \texttt{forbidEq\_smul} and then
    applying \texttt{downward\_forbidEq\_equal\_flags} on both sides
    yields the base-type equation
    $\texttt{downward}(c \cdot (A \cdot B))
       \,=_{F\text{-forbid}}\, \sum_j (c \cdot d_j)\,\texttt{downward}(F_j)$.
    The term is then moved to the right by
    \texttt{forbidLE\_rw\_left\_add\_right} together with
    \texttt{forbidLE\_move\_add\_left\_iff} (or, if it is the last
    term, the \texttt{\_rw\_left}/\texttt{\_move\_term\_\ldots}
    variants).
  \item If $T$ is a single flag term: no rewriting is needed; the
    tactic just applies
    \texttt{forbidLE\_move\_add\_left\_iff} (or
    \texttt{forbidLE\_move\_term\_left\_iff} for the last term) to
    move it to the right.
\end{itemize}

This process repeats until the left-hand side has nothing left to
process, so all \texttt{downward(\ldots)} and single-flag terms end
up on the right and the left side simplifies to a form ready to be
closed. Before the loop, a
\texttt{try simp only [downward\_add, add\_assoc]} distributes
\texttt{downward(a+b)} as \texttt{downward(a) + downward(b)} and
right-associates the sum.

The tactic is used in the three main proofs ― Mantel, $C_4$-Tur\'an,
and Erd\H{o}s pentagon ― always in the same pattern: unfold the
\texttt{flagQuadraticForm}s with \texttt{simp}, then call
\texttt{reduce\_downward\_flagmul} once. The Erd\H{o}s pentagon proof
in \texttt{LeanFlagAlgebras/API/ErdosPentagonAPI.lean} is the largest
of the three; three quadratic forms with $v_0, v_1, v_2$ of lengths
$8, 6, 5$ produce $8^2 + 6^2 + 5^2 = 125$ product terms, all
dispatched by a single call:

\begin{lstlisting}
theorem ErdosPentagon_flagAlgebra_API
    : C5.toFlagAlgebra <=[K3.toFinFlag] (24 / 625 : Real) * (1 : FlagAlgebra empty) := by
  ...
  simp [flagQuadraticForm, v_0, P_real, ratMatrixToReal, P, Fin.sum_univ_eight, add_assoc]
  simp [v_1, Q_real, ratMatrixToReal, Q, Fin.sum_univ_six, add_assoc]
  simp [v_2, R_real, ratMatrixToReal, R, Fin.sum_univ_five, add_assoc]
  reduce_downward_flagmul
  expand_one_at 5
  ...
\end{lstlisting}

The Mantel proof has one quadratic form of length $2$ ($4$ product
terms) and $C_4$-Tur\'an has two of lengths $4, 3$ ($25$ terms); in
both cases the corresponding \texttt{reduce\_downward\_flagmul} call
is again a single line.

\subsection{\texttt{prove\_flag\_expand\_with\_forbidden\_flag N} and
\texttt{prove\_flag\_mul\_with\_forbidden\_flag N}}

Both tactics are defined in \texttt{LeanFlagAlgebras/Logic/Tactic.lean}
and close \texttt{Assert}-level entailment goals of the form
\[
   F_{\text{forbidden}} =_\alpha 0 \;\vdash_\alpha\;
   \mathrm{LHS} =_\alpha \sum_j c_j\,F'_j,
\]
where the right-hand side is a linear combination of size-$N$ flags
that do \emph{not} contain $F_{\text{forbidden}}$. The two tactics
differ only in the shape of the left-hand side:

\begin{itemize}
  \item \texttt{prove\_flag\_expand\_with\_forbidden\_flag N} handles
    $\mathrm{LHS} = F$, a single flag: ``rewriting $F$ as a size-$N$
    expansion in which all forbidden terms vanish''.
  \item \texttt{prove\_flag\_mul\_with\_forbidden\_flag N} handles
    $\mathrm{LHS} = F_1 \cdot F_2$, a product of two flags: ``rewriting
    the product as a size-$N$ linear combination, again with the
    forbidden terms removed''.
\end{itemize}

A representative use from
\texttt{LeanFlagAlgebras/Logic/MantelTheorem.lean}, with
$F_{\text{forbidden}} = \mathrm{FlagAlgebra}\_3\_0\_0\_3$ (the size-$3$
$K_3$ flag) for \texttt{\_expand\_} and
$\mathrm{FlagAlgebra}\_3\_1\_0\_5$ for \texttt{\_mul\_}:

\begin{lstlisting}
example : FlagAlgebra_3_0_0_3 =a 0
    |-a FlagAlgebra_2_0_0_1
        =a (1 / 3 : Real) * FlagAlgebra_3_0_0_1
         + (2 / 3 : Real) * FlagAlgebra_3_0_0_2 := by
  prove_flag_expand_with_forbidden_flag 3

example : FlagAlgebra_3_1_0_5 =a 0
    |-a FlagAlgebra_2_1_0_0 * FlagAlgebra_2_1_0_0
        =a FlagAlgebra_3_1_0_2 + FlagAlgebra_3_1_0_0 := by
  prove_flag_mul_with_forbidden_flag 3
\end{lstlisting}

Internally both tactics follow the same script: parse
\texttt{FlagAlgebra\_n\_k\_m\_i} constants in the goal with
\texttt{parseFlagAlgebraIndices?}, combine the indices with the
argument \texttt{N} to synthesize the names
\texttt{flagSet\_N\_k\_m\_eq\_univ} and
\texttt{flagSet\_N\_k\_m\_val\_eq} (registered by the meta-commands),
then run a fixed sequence of \texttt{simp}/\texttt{rw} steps:
\texttt{intro $\varphi$ h}, expand the left-hand side as a sum over
size-$N$ flags (\texttt{unitVector\_quot\_eq\_sum} for the single-flag
case, \texttt{unitVector\_quot\_mul\_eq\_flagMul\_quot} for the
product case), unfold the resulting sum into an explicit list,
distribute the \texttt{PositiveHom}, substitute the hypothesis
\texttt{h} to drop the forbidden term, and finish with
\texttt{ring\_nf}. The \texttt{\_mul\_} variant additionally applies
\texttt{rw [mul\_comm]} first if the two factor indices are in the
wrong order, matching the ``smaller index on the left'' convention
used by \texttt{load\_forbid\_mul\_theorems}. Each of these
``$F$ vs.\ size-$N$ expansion'' and ``$F_1 \cdot F_2$ vs.\ size-$N$
expansion'' goals is closed by a single parameterized call instead
of the same 6--7-line script repeated by hand.

\subsection{\texttt{flag\_nonneg}}

A macro tactic defined in \texttt{LeanFlagAlgebras/API/Basic.lean}. It
closes goals of the form
\[
   0 \;\le\; c_1 \cdot \llbracket\mathrm{unitVector}\,F_1\rrbracket
   + c_2 \cdot \llbracket\mathrm{unitVector}\,F_2\rrbracket
   + \cdots + c_k \cdot \llbracket\mathrm{unitVector}\,F_k\rrbracket
   \qquad (c_i \ge 0),
\]
i.e.\ ``the right-hand side is a non-negative combination of unit
vectors''. This is the final step in all four main API proofs
(\texttt{API/MantelTheoremAPI.lean}, \texttt{API/C4TuranAPI.lean},
\texttt{API/K4freeP4.lean}, \texttt{API/ErdosPentagonAPI.lean}): the
earlier tactics (\texttt{reduce\_downward\_flagmul},
\texttt{expand\_one\_at}, \texttt{ac\_sort\_rhs\_pipeline}, \ldots)
move every nontrivial term to the right and reduce the left side to
$0$, and \texttt{apply forbidLE\_of\_le} drops the forbid qualifier;
each of the four proofs is concluded by

\begin{lstlisting}
  ...
  apply forbidLE_of_le
  flag_nonneg
\end{lstlisting}

It introduces an arbitrary positive homomorphism $\varphi$,
distributes it over the sum with \texttt{PositiveHom.map\_add},
splits leaf-by-leaf with \texttt{repeat apply add\_nonneg}, and
discharges each leaf with \texttt{positiveHom\_unitVector\_ge\_zero}.
\texttt{flag\_nonneg} bundles these four steps into a single macro.

\section{Algebraic normalization tactics}

After the domain-specific tactics resolve all products, the goal typically
reduces to an equality between two linear combinations sharing the same summands
in different orders. The \texttt{ac\_sort} family, defined in
\texttt{LeanFlagAlgebras/Utils/SortTactic.lean}, handles this final step.

An earlier attempt, \texttt{sort}, closed the sorted equality with
\texttt{ring\_nf} or \texttt{abel\_nf}, which proved too slow for the
tens-to-hundreds of summands typical in flag algebra goals. \texttt{ac\_sort}
keeps the sort but restricts the equality proof to add-AC only. The original
\texttt{sort} variants remain in the file but are unused in practice.

\texttt{ac\_sort} proceeds in four steps. \texttt{flattenAddTerms} flattens the
expression tree into a summand array $[t_1, \ldots, t_m]$ (e.g.\
\texttt{c~{\small\textbullet}~FlagAlgebra\_5\_0\_0\_17} or
\texttt{-~FlagAlgebra\_5\_0\_0\_3}). The sort key for each summand is obtained
by stripping the scalar with \texttt{getSmulArgs?} and reading the trailing
index of the base identifier with \texttt{parseTrailingNat?}.
\texttt{insertSortedBy} stable-sorts the array and \texttt{mkRightAssocAdd}
reassembles the result into $t_{\pi(1)} + \cdots + t_{\pi(m)}$. The auxiliary
equality is closed by
\texttt{first | ac\_rfl | simp [add\_assoc, add\_left\_comm, add\_comm]} and
rewritten into the goal via \texttt{conv}. Variants \texttt{ac\_sort\_lhs} /
\texttt{\_rhs} sort one side; \texttt{ac\_sort} sorts both.

The \texttt{ac\_sort\_pipeline} variants (\texttt{\_lhs\_pipeline} /
\texttt{\_rhs\_pipeline}) add pre- and post-processing around the sort. A
\texttt{simp only [neg\_add, neg\_neg, sub\_eq\_add\_neg,
$\leftarrow$~neg\_smul, add\_assoc, smul\_smul]} pass normalizes subtraction
and flattens nested scalar multiples before sorting; afterwards,
\texttt{simp only [$\leftarrow$~add\_assoc, $\leftarrow$~add\_smul]} and
\texttt{norm\_num} merge summands sharing the same base and evaluate
coefficients.

In the main proofs, \texttt{ac\_sort\_pipeline} (or its \texttt{\_rhs\_}
variant) closes the goal immediately after \texttt{reduce\_flagmul} or
\texttt{reduce\_downward\_flagmul}:

\begin{lstlisting}
  ...
  reduce_flagmul
  apply Forbid.forbidEq_of_eq
  ac_sort_pipeline
\end{lstlisting}

\section{Summary}

Each main API proof in this project follows the same skeleton. The
\texttt{load\_*} meta-commands run once at the module level, reading
externally generated JSON files and registering all flag definitions, density
values, and product expansion lemmas. Inside the proof body, \texttt{simp}
unfolds the quadratic form into an explicit sum of products;
\texttt{reduce\_downward\_flagmul} (or \texttt{reduce\_flagmul} for equality
goals) resolves every product into a linear combination of base flags;
\texttt{expand\_one\_at~n} rewrites the remaining $1$ as a weighted sum of
forbidden-free flags of size $n$; \texttt{ac\_sort\_rhs\_pipeline} normalizes
the right-hand side; and \texttt{apply~forbidLE\_of\_le} followed by
\texttt{flag\_nonneg} discharges the non-negativity condition. The Mantel,
$C_4$-Tur\'an, and Erd\H{o}s pentagon proofs all fit this template, differing
only in the size argument to \texttt{expand\_one\_at} and the dimensions of
the quadratic forms involved.

The key design decision is that definition and theorem names carry their
numeric indices --- \texttt{Flag\_n\_k\_m\_i}, \texttt{FlagAlgebra\_n\_k\_m\_i},
\texttt{flagMul\_<A>\_<B>}, and so on --- so each tactic can determine which
lemma to apply by inspecting the goal's \texttt{Expr} and synthesizing the
corresponding name, with no additional configuration at the call site. The
meta-commands that register these names and the tactics that consume them share
a single naming convention, which also makes the system easy to extend:
supporting a new forbidden graph or a larger host size amounts to generating a
new JSON file and adding one \texttt{load\_*} call, after which all existing
tactics continue to work without modification. The Mantel ($K_3$-free),
$C_4$-Tur\'an, and Erd\H{o}s pentagon ($K_3$-free) proofs already demonstrate
this: despite targeting different extremal problems, they share the same tactic
stack.

None of the tactics described here implement a sophisticated algorithm. Their
value is that the same boilerplate --- unfolding a flag expansion, resolving a
product, closing a non-negativity goal --- that would otherwise repeat across
every lemma is replaced by a single parameterized call. The steps that remain
visible in the proof, \texttt{reduce\_downward\_flagmul},
\texttt{expand\_one\_at}, and \texttt{flag\_nonneg}, correspond directly to the
mathematical operations on paper, which keeps the proofs concise and auditable.

\end{document}
